Krátké shrnutí: tato kapitola se věnuje výzvám hodnocení v době snadno dostupné generativní AI a nabízí praktický, krok‑za‑krokem přístup k redesignu zadání a hodnoticích postupů na úrovni předmětu i katedry. Cílem je pomoci garantům předmětů, lektorům a asistentům navrhnout hodnocení, které je pedagogicky smysluplné, transparentní vůči studentům a odolné vůči triviálnímu zneužití automatizovaných nástrojů.

Krátké opodstatnění přístupu:
- Namísto snahy „zakázat“ nástroje doporučuje materiál upřednostnit redesign úloh a pilotní nasazení tak, aby hodnocení měřilo procesy, dovednosti a reflektivní metody, nikoli pouze konečný text či kód (general background knowledge).  
- Pro pilotování a plánování strategie lze efektivně využít moderní enterprise agenty (např. Microsoft 365 Copilot) k rychlé tvorbě návrhů, plánů a analýz potřebných kroků; zkušenosti ukazují, že specializované agenty („Výzkumník“) dokážou výrazně zrychlit a zkvalitnit přípravu strategických dokumentů \cite{kb_ai_101_tipu_pdf_04e4a115_page_48_1}.  
- Pilotní nasazení v malém měřítku (seminář, laboratorní cvičení) umožní ověřit pedagogický přínos a doladit postupy před širším rozšířením — analogicky k doporučenému postupu při zavádění aktivit jako Minecraft Education nebo nástrojů pro procvičování matematiky \cite{kb_ai_101_tipu_pdf_04e4a115_page_15_2,kb_ai_101_tipu_pdf_04e4a115_page_8_1}.

Doporučený pracovní postup pro garanty a lektory (rychlý checklist)
\begin{enumerate}
  \item Zmapujte stávající hodnocení: identifikujte typy úloh (výsledkové vs. procesní), body rizika z hlediska automatizovaných nástrojů a které artefakty jsou prokazatelně unikátní (repo‑specifické, lokální data apod.). (general background knowledge)
  \item Rozhodněte o cílech redesignu: které kompetence mají být měřeny (proces, metakognice, aplikace), kde je vhodné povolit AI s povinnou dokumentací a kde je třeba časově/organizačně kontrolované prostředí. (general background knowledge)
  \item Navrhněte a spusťte malý pilot: vyberte jeden modul nebo seminář, připravte nové zadání/rubriky a metodu sběru důkazů (milníky, drafty, portfolia). Vyhodnoťte pilot a sepište lessons‑learned; inspiraci pro pilotní kroky najdete v praktických příkladech (Minecraft Education, nástroje pro „pokrok v matematice“) \cite{kb_ai_101_tipu_pdf_04e4a115_page_15_2,kb_ai_101_tipu_pdf_04e4a115_page_8_1}.
  \item Podpořte práci nástroji pro tvorbu a plánování: při přípravě sylabů, rubrik a checklistů lze využít Copilot‑typ agenta pro generování návrhů a konsolidaci podkladů; vždy ověřte generované výstupy lidskou kontrolou \cite{kb_ai_101_tipu_pdf_04e4a115_page_48_1}.
  \item Formalizujte změny do sylabů a komunikačních šablon: doplňte pole pro deklaraci použití AI, popište pravidla eskalace a spot‑checkové mechanismy pro kontrolu kvality. (general background knowledge)
\end{enumerate}

Co v kapitole najdete (stručný přehled)
- Šest principů AI‑aware assessment a jejich pedagogické důsledky;  
- Katalog úloh odolných vůči triviálnímu využití AI (ústní zkoušky, time‑boxed writings, reflexivní portfolia apod.);  
- Praktické rubriky, milestone‑based hodnocení a šablony pro transparentní deklaraci použití AI;  
- Nástroje pro formativní zpětnou vazbu a tipy pro technické obory;  
- Etické a GDPR aspekty hodnocení, důraz na dokumentaci zpracování a DPIA;  
- Šablony, checklisty a prompt‑knihovna pro rychlé nasazení.

Krátké doporučení na závěr: začněte malým pilotem, využijte dostupné nástroje pro rychlé návrhy a plánování, a výsledky pilotu převeďte do jasně komunikovaných změn v sylabu a rubrikách — tento cyklus „navrhni‑pilotu‑ověř‑škáluj“ je prioritou při adaptaci hodnocení na VUT. \cite{kb_ai_101_tipu_pdf_04e4a115_page_48_1,kb_ai_101_tipu_pdf_04e4a115_page_15_2,kb_ai_101_tipu_pdf_04e4a115_page_8_1}